\documentclass[10pt,twocolumn]{article}

% --- portable preamble: only packages shipped with any TeXLive ---
\usepackage[T1]{fontenc}
\usepackage[utf8]{inputenc}
\usepackage[margin=2cm]{geometry}
\usepackage{microtype}
\usepackage{booktabs}
\usepackage{xcolor}
\usepackage{listings}
\usepackage{amssymb}
\usepackage[hidelinks]{hyperref}

\lstset{
  basicstyle=\ttfamily\footnotesize,
  columns=fullflexible,
  keepspaces=true,
  breaklines=true,
  showstringspaces=false,
  frame=single,
  framesep=4pt,
  xleftmargin=2pt,
  aboveskip=6pt,
  belowskip=4pt,
}

\title{\textbf{Region Ownership: Unifying the Priority-Ceiling Protocol\\
and DMA Handshakes in a Bare-Metal Language}}

\author{Daniel Tralamazza\thanks{beemel was developed iteratively with AI
assistance; the design emerged from hardware bring-up rather than an up-front
specification. The author is responsible for all claims; the complete
development history and test suite are public.}\\
\small Independent\\
\small \texttt{tralamazza@gmail.com}}

\date{\textbf{Working draft} --- \today}

\begin{document}
\maketitle
\begin{center}\small\emph{Working draft, not peer-reviewed. The evaluation is
preliminary and being strengthened; the implementation is alpha and unproven.
Comments welcome.}\end{center}
\vspace{0.5em}

\begin{abstract}
A compiler for a microcontroller normally models a single bus master: the CPU.
But memory correctness on a modern MCU is a property of the whole system --- CPU
cores, DMA engines, caches, and the bus matrix that connects them. The resulting
bugs are silent: a buffer placed in memory the DMA engine cannot reach, a cache
line the DMA never sees, a descriptor armed before the CPU has finished writing
it. We describe BML, a language and compiler (beemel) for ARM Cortex-M firmware,
built around two ideas. First, every bus master is a first-class \emph{agent}
with its own reach, cache view, and permissions, and CPU/DMA memory sharing is
checked at compile time. Second --- our central claim --- the CPU
priority-ceiling protocol (mutual exclusion between interrupt and thread
contexts) and the asynchronous DMA release/reclaim handshake are \emph{one
concept, region ownership}, with the synchronization \emph{mechanism} derived
from the set of sharers --- the programmer declares the sharing and the ownership
window, not which exclusion protocol to emit. The model was
built failure-first: every check traces to a specific hardware bug we provoked
on a development board, and the result was validated back on silicon across
three DMA architectures (ST crossbar, RP2350 per-channel, nRF51 fixed-block). We
position the work against the priority-ceiling literature (which covers the CPU
half only) and the formal DMA-isolation and OS memory-model literature (which
treats DMA but never unifies it with CPU-side mutual exclusion). These are
compile-time analyses plus a test suite, not formal proofs; we are explicit
about the resulting trust boundary.
\end{abstract}

\section{Introduction}

The mental model behind most embedded toolchains is that a program runs on one
processor, and memory is a flat array that processor reads and writes. A
compiler optimizes against that model and a linker places symbols into it. The
model is wrong for any MCU with a DMA engine. A network controller, USB
controller, or timer-driven peripheral DMA is a second bus master: it issues its
own loads and stores, on its own clock, with its own view of memory --- a
different reachable address range over the bus matrix, and, on cached parts, a
view that does not snoop the CPU's data cache.

When the CPU and a DMA engine disagree about a region of memory, the failure is
typically silent. The program compiles and links. The board boots. And then it
never transmits a packet, or it transmits stale data, or it locks up
intermittently under load. The toolchain emitted no diagnostic because, in its
model, there was only ever one actor touching memory.

This paper describes BML, a language for Cortex-M firmware whose compiler,
beemel, models the second actor. Our contributions are:

\begin{itemize}
\item \textbf{The agent abstraction} (\S\ref{sec:agents}): a uniform treatment
of CPU cores, DMA engines, debug probes, and external firmware as
\emph{agents} --- bus masters with a declared reach, cache view, and
permissions --- with an admission schema that survived cross-vendor falsification
(no seventh question emerged across three DMA architectures).

\item \textbf{Region ownership} (\S\ref{sec:ownership}): the central
contribution. We show that the priority-ceiling protocol (CPU/CPU mutual
exclusion) and the DMA release/reclaim handshake (CPU/DMA mutual exclusion)
impose the \emph{same ownership obligation}, and that a compiler can
\emph{derive} which exclusion mechanism to emit from the set of agents that share
a region --- so the two compose automatically when an agent and a CPU context
share one buffer. To our knowledge no prior source language derives the two from
a single sharer-set abstraction (\S\ref{sec:related}).

\item \textbf{Derived whole-system checks} (\S\ref{sec:checks}): bus-matrix
reachability, cache-coherence enforcement by generated MPU regions, transfer
extent, and volatility-by-provenance, each derived from placement and use rather
than annotated.

\item \textbf{An empirical, failure-driven evaluation}
(\S\ref{sec:eval}): every check is traced to a concrete bug provoked on
hardware, and the model is validated on three boards. We report the
abstract-interpretation encodings that were required to discharge the DMA
obligations in a non-relational domain.
\end{itemize}

We are deliberate about what this is not. BML has no mechanized semantics and no
soundness theorem; its guarantees are compile-time analyses backed by an
end-to-end test suite and optional static verification. \S\ref{sec:limits}
states the trust boundary precisely. We argue the contribution is a \emph{model
and its validation}, in the tradition of language-design experience reports.
Figure~\ref{fig:example} shows a complete BML fragment up front; the rest of the
paper explains what each line checks.

\begin{figure*}[t]
\begin{lstlisting}
import stm32h723.mdma;

owns MDMA;                                 // exclusive control of the engine

var SRC: [u8; 64] in dma_shared;           // MDMA reads this region
var DST: [u8; 64] @shared in tcm_scratch;  // MDMA writes; thread + ISR also read
var RESULT: u32;                           // CPU-visible result

fn copy() {
    var i: u32 = 0;                        // fill before release: index-WRITE is legal
    while i < 64 { SRC[i] = i as u8; i = i + 1; }

    MDMA.MDMA_C0SAR = &SRC as u32;         // release: hand the engine the source ...
    MDMA.MDMA_C0DAR = &DST as u32;         // ... and the destination
    MDMA.MDMA_C0CR.EN = true;              // arm

    claim DST {                            // CPU window: masks the ISR co-sharer
        if MDMA.MDMA_C0ISR.CTCIF0 {        // observe the engine's completion flag
            const v = reclaim(DST);        // only now is a bounds-checked view sound
            RESULT = v[0] as u32;          // (calls are forbidden inside a claim)
        }
    }
}
\end{lstlisting}
\caption{A complete BML fragment, condensed from the H723 MDMA example.
\texttt{DST} is shared by the thread, an ISR, and the MDMA engine, so its carrier
is \texttt{Shared(AgentShared([u8;64]))} and every access outside the window is
rejected: dropping the \lstinline!claim! makes \lstinline!reclaim! an error (the
ISR could race), dropping the flag test is a missing-handshake error, and an
index-read of \texttt{DST} elsewhere is rejected because the engine may still own
it. The \lstinline!claim!/\lstinline!reclaim! nesting --- one CPU window, one
agent handshake --- is the unification of \S\ref{sec:ownership}, emitted from the
fact that a CPU context and a DMA agent share \texttt{DST}.}
\label{fig:example}
\end{figure*}

\section{Motivation: a CPU-only compiler is blind}
\label{sec:motivation}

The model in this paper was not designed up front. It was forced out by a
sequence of failures on a NUCLEO-H723ZG (Cortex-M7) while bringing up an
Ethernet driver, and refined against two further boards. The founding failures,
all observed:

\paragraph{Unreachable placement.} Moving the transmit descriptor ring into
DTCM (tightly-coupled memory) compiled, linked, and produced a board that never
transmitted. The Ethernet DMA cannot address DTCM over the bus matrix; nothing
in the toolchain knew the bus topology, so nothing objected.

\paragraph{Unchecked cache discipline.} Synchronizing the descriptor publish
with a `dmb` barrier alone was sound only because the D-cache happened to be
disabled. Enabling the cache would have broken receive silently --- the DMA
writes a descriptor the CPU never sees because the CPU reads a cached copy.

\paragraph{A posted-write completion hazard.} Arming the DMA is a posted
(buffered) write on a Device-memory register. One left in flight while the bus
remained busy was an intermittent source of imprecise bus faults --- present or
absent depending on instruction scheduling.

\paragraph{An optimizer-hoisted ownership poll.} A plain load of a descriptor's
ownership bit, in a spin loop waiting for the DMA to finish, was hoisted by the
optimizer out of the loop into an infinite branch-to-self. The CPU could not see
that a concurrent writer (the DMA) existed.

Each failure is a disagreement between two bus masters: about reachability,
about visibility, about completion ordering. A compiler that models only the CPU
has no way to state any of them, let alone check them. The rest of the paper is
the vocabulary we added to state them and the checks that follow.

\section{The agent model}
\label{sec:agents}

\textbf{An agent is anything that touches memory on its own initiative.} The
admission test is three questions, all of which must hold: (1) does it initiate
memory accesses itself, rather than only decode them? (2) does it act
concurrently, on its own clock, once set up? (3) does it have its own view of
memory --- reach, caching, permissions? A timer that triggers a peripheral DMA
is an agent; the peripheral's own register file is not. The RP2350 PIO executes
instructions but has no memory reach, so it fails question (3) and is not an
agent.

\begin{table}[t]
\centering
\footnotesize
\begin{tabular}{@{}lll@{}}
\toprule
kind & example & accesses answered for by \\
\midrule
\texttt{cpu}      & Cortex-M core           & the compiler (every IR access) \\
\texttt{dma}      & Ethernet DMA, EasyDMA   & owning module, at handoffs \\
\texttt{debug}    & SWD probe               & host side; concurrency-inert \\
\texttt{external} & other-vendor firmware   & nobody; channels only \\
\bottomrule
\end{tabular}
\caption{Agent kinds. A \texttt{cpu} agent's accesses are visible in the
compiler's IR and checked per access; a \texttt{dma} agent's accesses are
invisible, so the compiler sees only the address flowing into a register.}
\label{tab:agents}
\end{table}

The asymmetry in Table~\ref{tab:agents} is the crux. For a \texttt{cpu} agent
the compiler emits every load and store, so its rules are ordinary per-access
checks. For a \texttt{dma} agent the accesses are opaque --- the compiler sees
only an address written into a controller register. A \emph{handoff} is that
register write, and it is where the compiler recovers, for an opaque agent, the
visibility it gets for free on the CPU.

The agent contract is, so far, a closed schema: every key in an agent declaration
answers one of six questions about a bus master --- may it run, where can it
touch, which buffer, how much, when is it done, what code does it run. We arrived at the six
by reviewing and instantiating three structurally different DMA architectures:
ST's central crossbar with software-selected master ports (H723 Ethernet and
MDMA), Raspberry Pi's per-channel controller (RP2350), and Nordic's
per-peripheral fixed-block engines (nRF51, where every capable peripheral is its
own bus master). No seventh question emerged across the three; the schema is
closed only modulo the acknowledged future refinements of \S\ref{sec:limits}
(e.g. scatter-gather descriptor chains), not proven minimal.

\section{Region ownership}
\label{sec:ownership}

A \emph{region} names a block of memory and the set of agents that share it.
This section is the core claim: a CPU sharer and a DMA sharer impose the same
ownership \emph{obligation} --- exclusive ownership with observed transfer --- and
the compiler derives the (distinct) exclusion mechanism from the sharer set.

\subsection{Two masters, one ownership}

Consider a buffer shared between the CPU and one other agent. Correct access
requires that at any moment at most one of them treats the buffer as live, and
that ownership transfers are observed by both sides. This is mutual exclusion,
and it has two textbook instantiations that are normally studied in separate
literatures:

\begin{itemize}
\item If the other sharer is \textbf{another CPU context} (an ISR sharing a
buffer with a thread, or a higher-priority ISR), exclusion is achieved by
raising the processor's priority above every context that uses the buffer --- the
\emph{priority-ceiling protocol}~\cite{baker1991}. The owner holds the buffer
for the duration of a masked critical section.

\item If the other sharer is a \textbf{DMA engine}, exclusion is achieved by an
\emph{asynchronous handshake}: the CPU writes the buffer's address into a
controller register (\emph{release}), the engine works on it, and the CPU must
observe a completion signal before touching the buffer again (\emph{reclaim}).
\end{itemize}

We claim these share one \emph{obligation} --- region ownership: at most one live
owner, transfers observed by both. The two \emph{mechanisms} remain genuinely
different --- ceiling masking is preventive and synchronous (the contending
context cannot run), the DMA handshake is cooperative and asynchronous (release
now, observe completion later) --- and which one the compiler emits is a function
of \emph{who} the co-sharer is. The contribution is not that the mechanisms are
identical, but that one obligation, derived from the sharer set, subsumes both
and makes their composition automatic.

\subsection{The CPU sharer: the claim window}

When a region is shared between CPU contexts, BML attaches a \emph{ceiling} to
the shared static, computed from the set of contexts that access it. A
\lstinline!claim X { ... }! block opens an ownership window: on ARMv7-M it raises
\texttt{BASEPRI} to the static's ceiling, so unrelated higher-priority
interrupts keep running while every \emph{contending} context is masked; on
ARMv6-M, which lacks \texttt{BASEPRI}, it falls back to a global
\texttt{cpsid}/\texttt{cpsie} pair. Inside the window the static is at its
inner type --- ordinary reads and writes are legal --- and the per-access critical
sections the compiler would otherwise insert are suppressed, because the window
already provides exclusion.

\subsection{The DMA sharer: release and reclaim}

When a region is shared with a DMA agent, an array placed in it is given a
move-typed carrier at name resolution. \emph{Release} is the handoff register
write. \emph{Reclaim} is an explicit \lstinline!reclaim(BUF)! that yields a
bounds-checked view and is legal only once the channel's completion flag has been
observed. The completion guard is checked by lexical containment: the
\lstinline!reclaim! must sit within the guard span of an observation of the
channel's own \texttt{completes\_by} flag (a conservative, non-flow-sensitive
check) (Table~\ref{tab:target}). A plain view over
move-typed memory, without the handshake, is rejected and the diagnostic points
at \lstinline!reclaim!.

\subsection{Deriving the mechanism, and composing both}

The programmer does not choose between \lstinline!claim! and
\lstinline!reclaim!. The compiler reads the region's sharer set and emits the
mechanism: a masked window for a CPU co-sharer, a release/reclaim handshake for a
DMA co-sharer. When both are present the carriers nest, and unguarded access
becomes unrepresentable in well-typed code: under a static that is both
ceiling-shared and DMA-shared, \emph{every} access --- reads included --- is
rejected outside a \lstinline!claim!, so the masked window is required by
construction and the handshake composes inside it with no additional rule:

\begin{lstlisting}
claim X {            // CPU co-sharers masked here
    if done() {      // DMA co-sharer's completion observed
        reclaim(X)   // only now is a view over X sound
    }
}
\end{lstlisting}

The same observation runs the other way at the \texttt{verify} layer
(\S\ref{sec:verify}): a read of the claimed static \emph{inside} the window is
discharged as stable by IKOS (the mask excludes the only writer), while the
identical read outside the window is not. Ownership is the single fact that both the type
checker and the abstract interpreter consult.

\section{Whole-system checks derived from the model}
\label{sec:checks}

The ownership story rests on a layer of checks that connect placement to
physics. They are organized so that the truth lives in exactly one place: chip
\emph{physics} in a target file, written once per part; project \emph{policy} as
regions; and \emph{software binding} in the source. The governing principle is
that the program declares only what use cannot reveal --- the bus matrix and
exclusivity claims --- and the compiler derives the rest.

\begin{table}[t]
\centering
\footnotesize
\begin{lstlisting}
[mem.dma_pool]
base = 0x30007000
size = 4K
cacheable = false       # -> generated MPU region

[agent.mdma]
kind        = dma
reach       = dma_pool, dtcm   # checked vs. bus
bus         = axi: 0x08000000..0x08100000,
              ahbs: 0x00000000..0x00040000
enabled_by  = RCC.C1_AHB3ENR.MDMAEN

[agent.mdma.ch0]
completes_by = MDMA.MDMA_C0ISR.CTCIF0
extent       = MDMA.MDMA_C0BNDTR.BNDT
handoff      = MDMA.MDMA_C0SAR
               port_by MDMA.MDMA_C0TBR.SBUS ahbs

[region.dma_shared]
mem    = dma_pool
agents = mdma           # CPU always implicit
\end{lstlisting}
\caption{A target-file fragment (Layer 1, physics). No key takes a raw address:
every register is an SVD path, so a stray literal fails at resolution rather than
silently configuring hardware.}
\label{tab:target}
\end{table}

\paragraph{Bus-matrix reachability.} A region's memory must lie within every
listed agent's reach, and \texttt{reach} is itself cross-checked against the
transcribed \texttt{bus} windows. The DTCM placement bug from
\S\ref{sec:motivation} fails at target load, before any source is compiled.

\paragraph{Cache coherence as generated protection.} A cacheable block shared by
a cached CPU and a non-snooping DMA agent is rejected. Declaring it
\lstinline!cacheable = false! does not merely annotate: it \emph{generates} a
non-cacheable MPU region in the reset handler (PMSAv7 on Cortex-M4/M7, PMSAv8 on
M33). Alignment of a static in such a region is floored to the CPU's cache-line
size --- the alignment requirement was always a property of the cache, never of
the DMA engine.

\paragraph{Transfer extent.} The \texttt{extent} field names the controller's
transfer-count register; at the \texttt{verify} layer, arming the channel emits
an obligation that \texttt{count} $\times$ \texttt{unit} does not exceed the
delivered buffer's capacity.

\paragraph{Volatility by provenance.} BML has no \texttt{volatile} keyword.
Volatility is a property of where data lives: a pointer derived from the address
of a DMA-shared static points at memory with a concurrent writer the optimizer
cannot see, so every load and store through such a pointer is lowered as
volatile. The taint is tracked syntactically from the address-of site to a
fixpoint, and an escape check (the pointer may not leave the deriving function)
prevents it from being silently lost; an address laundered through integer
arithmetic escapes the syntactic taint and is a recorded blind spot
(\S\ref{sec:limits}). This directly addresses the optimizer-hoisting failure of
\S\ref{sec:motivation}, and is a concrete response to the long-standing
observation that \texttt{volatile} is fragile and miscompiled~\cite{eide2008}.
Crucially, a view obtained \emph{inside} a justified ownership window stays
non-volatile: the agent is excluded there, so the proof restores the
optimization that the unguarded boundary forbids.

\section{Optional verification}
\label{sec:verify}

\subsection{The pipeline}

\texttt{bml verify} lowers the program to LLVM IR and runs
IKOS~\cite{ikos}, an abstract-interpretation analyzer, for buffer overflows,
null dereferences, division by zero, integer overflow, and user assertions. The
regions model adds obligations on top: handoff provenance and reach, in-memory
descriptor pointer placement, and transfer extents.

\subsection{Encodings the domain forced}

Discharging the DMA obligations in IKOS's default non-relational
interval-with-congruence domain required specific encodings, which we report
because they are reusable lessons for anyone driving abstract interpretation at
this layer:

\begin{itemize}
\item The provenance \texttt{assume} must sit at the address-of site. Two integer
casts of the same global do not unify across function boundaries, so an assume
placed at the use site proves nothing.
\item Capacity is encoded as a \emph{constant} compared against a count interval.
A base/limit \emph{pair} of shadow variables is unprovable in a non-relational
domain; a constant-versus-interval comparison is trivial in it.
\item Alignment is carried by exact addresses, not congruences. The domain does
no backward congruence narrowing, and because the compiler owns the linker
script, the exact address is available and sufficient.
\item Claim-aware havoc. Reads of a \texttt{@shared} static are normally
invalidated before each load (a higher-priority writer may have run); reads of
the \emph{claimed} static inside its window are not, because the mask provides
exactly that stability. This is the ownership fact of \S\ref{sec:ownership} seen
by the analyzer.
\end{itemize}

\section{Evaluation}
\label{sec:eval}

We offer three kinds of evidence: that BML's checks reject hazards
\emph{independently documented} for chips (\S\ref{sec:eval-catalog}) and
\emph{third-party drivers} (\S\ref{sec:eval-thirdparty}) it was not co-designed
with; a controlled counterfactual reproduced on hardware
(\S\ref{sec:eval-counterfactual}); and end-to-end bring-ups on three boards
(\S\ref{sec:eval-boards}). The first two directly address the central threat to
an experience paper --- that the checks were exercised only on the programs they
were built against.

\subsection{Documented hazards caught at compile time}
\label{sec:eval-catalog}

\begin{table*}[t]
\centering\footnotesize
\begin{tabular}{@{}p{3.0cm}p{2.7cm}p{3.4cm}p{5.9cm}@{}}
\toprule
Documented hazard (chip) & Attested by & Caught today by & BML, at compile time \\
\midrule
DMA buffer in DTCM (STM32H7) & AN4839~\cite{an4839}; ST community & nothing: board silently never transmits & rejected at target load: \emph{``region in mem \texttt{dtcm}, which agent \texttt{eth} cannot reach''} \\
DMA buffer in CCM RAM (STM32F4) & ST community; NuttX \texttt{CCMEXCLUDE} & runtime AHB / hard fault & rejected: \emph{``mem \texttt{ccm} \ldots agent \texttt{dma2} cannot reach''} \\
EasyDMA buffer in flash (nRF52) & Nordic spec~\cite{easydma}; \texttt{nrf-hal} \#37 & hard fault or silent; embassy's run-time \texttt{BufferNotInRAM} & rejected: \emph{``mem \texttt{flash} \ldots agent \texttt{easydma} cannot reach''} \\
Cacheable DMA buffer, D-cache on (STM32H7, SAME70) & AN4839~\cite{an4839}; Zephyr \#36471~\cite{zephyrdma} & silent corruption (``works with cache off'') & rejected: \emph{``cache views diverge''}; the fix \texttt{cacheable=false} auto-generates the non-cacheable MPU region \\
\bottomrule
\end{tabular}
\caption{Independently-documented, multi-vendor DMA hazards, each reconstructed
in BML and rejected by the compiler. One Layer-2 reach check (rows 1--3) catches
the same bug class across three bus architectures; the cache check (row 4) both
rejects the program and generates the vendor-prescribed fix. Reconstructions and
verbatim compiler output are in the artifact.}
\label{tab:hazards}
\end{table*}

Cortex-M firmware has a small family of notorious DMA bugs that vendors document
in application notes and that recur for years in forums and issue trackers. We
took four, spanning three vendors, and reconstructed each minimally in BML --- the
documented memory-sharing structure, not the original C, which BML cannot
ingest. In every case the compiler rejects the program;
Table~\ref{tab:hazards} summarizes, and the moved-to-reachable-memory
counterpart of each builds clean.

Two observations. First, a single check --- the Layer-2 reach test --- catches the
same bug class (a buffer the engine cannot address) across three structurally
different bus arrangements (ST's DTCM, ST's CCM, Nordic's flash exclusion): the
agent/reach model's generality made concrete rather than asserted. Second, the
``caught today by'' column is the point. These hazards are presently caught at
\emph{runtime} (a hard fault), \emph{not at all} (silent corruption), or by
\emph{manual} configuration one can omit; BML moves them to compile time, and for
the cache hazard it \emph{generates} the non-cacheable MPU region that ST's
examples write by hand.

\paragraph{Honest scope.} These are reconstructions, not analyses of the upstream
drivers: we re-expressed each documented hazard's placement and sharing in BML
and confirmed the diagnostic by running the compiler. The reach test is only as
good as the transcribed bus matrix, which lives in the shipped chip target
(\S\ref{sec:checks}); a mis-transcribed reach window would mask the bug. The model
also does not capture every subtlety --- the Cortex-M7 write-through-cache erratum
and speculative-read concerns of AN4839~\cite{an4839} are outside it. The claim is
narrower and concrete: the common, high-frequency form of each hazard becomes a
compile error.

\subsection{Third-party drivers: the same constraint, by hand}
\label{sec:eval-thirdparty}

The catalog cites vendor application notes; the same hazards recur in
\emph{third-party} drivers BML never saw, where named maintainers hit them and
added a guard --- weaker and later than a compile error. Two cases, each
reconstructed in BML and anchored to the project's own issue:

\paragraph{\texttt{stm32-rs/stm32-eth} (Rust), issue \#16~\cite{stm32eth}.} The
driver's maintainers documented that descriptors and buffers in a
D-cache-enabled region ``would need cleaning \ldots and invalidating,'' and
proposed to \emph{require} that they not be placed in such a region.
Reconstructing a cacheable TX descriptor ring handed to the ETH DMA, BML rejects
it at target load (\emph{``region \texttt{eth\_ring} \ldots cache views
diverge''}) and, with \texttt{cacheable = false}, generates the non-cacheable
MPU region that otherwise has to be arranged by hand (a dedicated linker section
plus a compile-time placement assertion --- the standard defensive pattern).

\paragraph{\texttt{nrf-rs/nrf-hal} and embassy (Rust), issue \#37~\cite{nrfhal}.}
``EasyDMA doesn't read from Flash'': a flash-resident \texttt{\&'static} buffer
handed to a UARTE transmit failed \emph{silently}. The fix was a \emph{run-time}
check returning \texttt{BufferNotInRAM} (embassy additionally offers a variant
that copies the data into a RAM scratch buffer). Reconstructing the same
delivery, BML rejects it at target load (\emph{``region \texttt{tx\_buf} in mem
\texttt{flash} \ldots agent \texttt{easydma} cannot reach''}) --- the same
physical fault, moved from run time to compile time, with no copy.

The pattern is the point: independent projects converged on BML's constraint and
enforced it the only ways a library can --- a runtime check, a defensive copy, or
a documented assertion --- where BML makes it a property the compiler checks once.
The honest scope of \S\ref{sec:eval-catalog} carries over: these reconstruct the
documented memory-sharing structure, not the upstream source, which BML cannot
ingest.

\subsection{Falsified both ways}
\label{sec:eval-counterfactual}

The volatility-by-provenance check of \S\ref{sec:checks} is a controlled
counterfactual on hardware. Under a plain (non-volatile) load, the Ethernet
controller's ownership-bit poll froze at transmit frame~5 --- the optimizer had
hoisted the load out of the spin loop into a branch-to-self. Under the
agent-derived volatile lowering, the identical program runs indefinitely. The bug
and its fix were each reproduced on the board, not only in the IR.

\subsection{Hardware bring-ups}
\label{sec:eval-boards}

Three boards span the three DMA architectures of \S\ref{sec:agents}. These are
existence proofs that the model is self-consistent and runs end to end; they do
not, on their own, isolate the unification claim (\S\ref{sec:limits}).

\paragraph{NUCLEO-H723ZG (Cortex-M7).} Ethernet transmit and receive with typed
descriptors, D-cache \emph{enabled} under the generated MPU region. An MDMA copy
into DTCM routes through the software-selected bus port. The full claim/reclaim
composition holds an exact cross-context invariant ($\mathrm{LOG\_SUM} = 4 \times
\mathrm{TICKS} - 10$) over a long run. Descriptor extents are proven at the
\texttt{verify} layer and confirmed at run time by the DMA's write-back of the
transferred length.

\paragraph{Pico 2 W (dual Cortex-M33).} First-flash boot via the Arm
\texttt{IMAGE\_DEF} block; the generated PMSAv8 MPU registers read back exactly
as emitted; the core1 launch handshake; and a cross-core \lstinline!claim!
backed by a hardware spinlock sustaining tens of millions of contended ownership
windows with no observed in-window invariant violation.

\paragraph{micro:bit v1 (nRF51, Cortex-M0).} An AES-ECB known-answer test driven
end to end through the full chain --- handoff, fixed-size extent, guarded reclaim
--- passing on first flash, with the ARMv6-M word-composed interrupt-priority path
exercised.

\section{Related work}
\label{sec:related}

The components of our model exist in prior work; the unification does not.

\paragraph{CPU mutual exclusion.} The priority-ceiling and Stack Resource
Policy~\cite{baker1991} are the CPU-side half of our ownership story, and they
are well established: Ada's Ravenscar profile~\cite{ravenscar} implements ceiling
locking, and RTIC~\cite{rtic} brings it to embedded Rust with
\texttt{BASEPRI}-based critical sections computed from the set of sharing tasks.
None of this work extends the protocol to a non-CPU sharer or an asynchronous
completion handshake.

\paragraph{DMA as a principal with a memory view.} The closest prior art to our
agent abstraction is the ETH Zurich decoding-net / least-privilege memory
model~\cite{achermann2019}, which treats DMA engines and cores as first-class
subjects in a network of address spaces and even pre-computes fixed-topology
translations at compile time. It is, however, an operating-system
memory-management model enforced by a kernel reference monitor through runtime
capabilities (implemented in Barrelfish, with a sketch for Linux); its reachability
machinery \emph{enables} an access by finding a translation path, assuming
reconfigurable MMU/IOMMU hardware. BML targets bare metal with no such hardware:
reach is a fixed property of the bus matrix, and an unreachable placement is a
compile error with no path to route around it. The decoding-net work models DMA
as a principal but never connects it to CPU-side mutual exclusion.

\paragraph{DMA buffer handoff.} Embedded Rust encodes handoff as affine move
semantics: a \texttt{Transfer} type consumes the buffer by value and returns it
only after a busy-wait on completion~\cite{embedonomicon}. This is a library
idiom rather than a language-level model; it does not check transfer length
against buffer capacity, and reachability is assumed. Crucially, RTIC and
embassy~\cite{embassy} already provide both ceiling-based locking and typestate
DMA ownership in one framework --- but as two distinct, programmer-chosen
mechanisms: neither derives the choice from the sharer set, and neither forces
the DMA handshake to nest inside the ceiling window when a buffer is shared both
ways. That single derivation, and the composition it makes automatic, is our
delta over them. Per-driver functional proofs, such as the Pancake verification
of an i.MX Ethernet driver~\cite{pancake}, verify one descriptor ring in program
logic, single-threaded, with no ceiling protocol.

\paragraph{Formal DMA isolation.} Haglund and Guanciale~\cite{haglund2022} give
a HOL4 framework that proves the conditions under which a DMA controller confines
its accesses to permitted memory regions --- a security-policy property,
discharged by interactive theorem proving for a driver or runtime monitor. The
framework models neither caches nor CPU-side concurrency, and it verifies a
runtime configuration rather than checking a source program.

\paragraph{Ownership types and verified DMA drivers.} Refined ownership types for
C --- RefinedC~\cite{refinedc} and CN~\cite{cn} --- formalize exactly the
ownership/capacity discipline our move-typed carriers and extent obligations
approximate, but as a general separation-logic framework, not a bus-master model
with reach, caches, or a ceiling protocol. Cogent~\cite{cogent} brings linear
types to verified systems code (file systems, not DMA paths). In the seL4 line, a
verified ZynqMP DMA driver in concurrent separation logic~\cite{seldma} treats
device and driver as concurrent agents over shared state --- the closest framing
to ours --- but verifies one driver per proof and does not unify the device
handshake with CPU-context exclusion.

\paragraph{Static analysis and compiler-mediated DMA isolation.} The closest
prior statement of our founding observation --- that one cannot verify a driver
without modeling the asynchronous device --- is Monniaux~\cite{monniaux2007}, who
verified an industrial USB OHCI driver by abstract interpretation, composed
asynchronously with a hand-written model of the DMA-capable controller, to prove
absence of run-time errors of exactly the kind where ``incorrect programming of
the device'' leads to ``memory zones being erased.'' This is general-purpose
abstract interpretation of a single driver-plus-device model, not a language in
which bus masters are first-class agents: it does not derive reachability, model
caches, or relate the device handshake to CPU-side mutual exclusion. PISTIS
\cite{pistis2022} isolates DMA on low-end MCUs without an IOMMU using a small
trust anchor together with a compiler pass that inserts mediation around accesses
to DMA-controller registers; enforcement is at run time by the monitor, not a
compile-time property of the program text.

\paragraph{Safe-language operating systems.} Singularity~\cite{singularity} made
device interfaces type-safe but explicitly placed DMA outside its safe model,
identifying it as the one inherently unsafe part of the driver/device interface.
Ivory and Tower~\cite{ivory} give a memory-safe embedded DSL whose ``regions''
are allocation/effect regions in the Tofte--Talpin sense, not hardware-shared
memory; Tock~\cite{tock} isolates untrusted capsules and processes with the MPU
at run time. None unifies a CPU ceiling protocol with a DMA handshake.

\paragraph{Summary.} Across the priority-ceiling, OS memory-model, formal
DMA-isolation, static-analysis, and safe-language-OS literatures, the components
recur but the unification does not: we found no source language that joins
CPU-side mutual exclusion and the DMA release/reclaim handshake into one
ownership abstraction, derived from the sharer set and checked at compile time.
That unification is the contribution.

\section{Limitations and threats to validity}
\label{sec:limits}

\paragraph{Checks, not proofs.} BML's compile-time analyses, its end-to-end test
suite (programs are executed under QEMU), and its optional IKOS verification
reduce a class of bugs. They are not a soundness guarantee. The compiler has no
mechanized semantics and has not been audited; it can have bugs of its own.

\paragraph{Recorded blind spots.} We maintain an explicit register of what is
trusted or lexically invisible. Among the entries: an address laundered through
integer arithmetic defeats the syntactic volatile taint; system-exception
priorities (SysTick, PendSV) are unmodeled; a completion flag cleared through a
\emph{separate} register is invisible to the straight-line staleness check; and a
fact carried across a loop back-edge is outside the lexical windows --- which
today yields a concrete false positive: a correct continuous-copy DMA loop whose
completion flag is cleared through a separate register (the STM32 IFCR idiom) is
rejected. Each is documented where it was discovered rather than elided.

\paragraph{Scope.} The implementation targets 32-bit ARM Cortex-M only. The
cross-vendor falsification covers three DMA architectures, which is evidence for
the schema's generality but not a proof of it; a controller with a structurally
new ``question'' would extend the contract.

\paragraph{Evaluation.} The evaluation is a set of hardware bring-ups, not a
controlled bug-finding study against a corpus of existing firmware. The strongest
threat to the central claim is that the model's value has been demonstrated on
the programs it was co-designed with; a study applying the checks to independent
third-party drivers is future work and would materially strengthen the case.

\section{Conclusion}

A microcontroller is a small distributed system: several bus masters, several
views of memory, no shared clock. Treating it as a single processor with a flat
memory is the source of a class of silent firmware bugs. BML treats every bus
master as an agent and, more pointedly, observes that the CPU and the DMA engine
want the \emph{same} thing from a shared buffer --- exclusive ownership, with
observed transfers --- and differ only in the mechanism, which the compiler can
derive from who is sharing. Pushed against real silicon across three DMA
architectures, the model suggests that whole-system memory correctness on an MCU
is a checkable property and that one ownership obligation spans both masters. The
implementation is alpha and the model unproven; \S\ref{sec:limits} states the
trust boundary.

\paragraph{Availability.} beemel, the BML language specification, the example
boards, and the verification toolchain are open source under Apache-2.0 at
\url{https://github.com/tralamazza/beemel}.

\begin{thebibliography}{99}
\footnotesize

\bibitem{baker1991} T.~P. Baker.
\newblock Stack-based scheduling of realtime processes.
\newblock \emph{Real-Time Systems}, 3(1):67--99, 1991.

\bibitem{rtic} J.~Aparicio, P.~Lindgren, et al.
\newblock RTIC: Real-Time Interrupt-driven Concurrency.
\newblock \url{https://rtic.rs}.

\bibitem{eide2008} E.~Eide and J.~Regehr.
\newblock Volatiles are miscompiled, and what to do about it.
\newblock In \emph{Proc. EMSOFT}, 2008.

\bibitem{achermann2019} R.~Achermann, N.~Hossle, L.~Humbel, D.~Schwyn, D.~Cock,
and T.~Roscoe.
\newblock A least-privilege memory protection model for modern hardware.
\newblock arXiv:1908.08707, 2019.

\bibitem{embedonomicon} The Embedonomicon: A no-std DMA API.
\newblock Rust Embedded Working Group.
\newblock \url{https://docs.rust-embedded.org/embedonomicon/dma.html}.

\bibitem{pancake} Verifying device drivers with Pancake.
\newblock arXiv:2501.08249, 2025.

\bibitem{haglund2022} J.~Haglund and R.~Guanciale.
\newblock Formally verified isolation of DMA.
\newblock In \emph{Proc. FMCAD}, 2022.

\bibitem{monniaux2007} D.~Monniaux.
\newblock Verification of device drivers and intelligent controllers: a case
study.
\newblock In \emph{Proc. EMSOFT}, pages 30--36, 2007.

\bibitem{pistis2022} M.~Grisafi, M.~Ammar, M.~Roveri, and B.~Crispo.
\newblock PISTIS: Trusted computing for low-end embedded systems.
\newblock In \emph{Proc. USENIX Security Symposium}, 2022.

\bibitem{singularity} G.~Hunt and J.~Larus.
\newblock Singularity: Rethinking the software stack.
\newblock \emph{ACM SIGOPS Operating Systems Review}, 41(2), 2007.

\bibitem{ivory} P.~C. Hickey, L.~Pike, T.~Elliott, J.~Bielman, and J.~Launchbury.
\newblock Building embedded systems with embedded DSLs.
\newblock In \emph{Proc. ICFP}, 2014.

\bibitem{tock} A.~Levy, B.~Campbell, B.~Ghena, et al.
\newblock Multiprogramming a 64\,kB computer safely and efficiently.
\newblock In \emph{Proc. SOSP}, 2017.

\bibitem{ikos} G.~Brat, J.~A. Navas, N.~Shi, and A.~Venet.
\newblock IKOS: A framework for static analysis based on abstract interpretation.
\newblock In \emph{Proc. SEFM}, 2014.

\bibitem{refinedc} M.~Sammler, R.~Lepigre, R.~Krebbers, K.~Memarian, D.~Dreyer,
and D.~Garg.
\newblock RefinedC: automating the foundational verification of C code with
refined ownership types.
\newblock In \emph{Proc. PLDI}, 2021.

\bibitem{cn} C.~Pulte, D.~C. Makwana, T.~Sewell, K.~Memarian, P.~Sewell, and
N.~Krishnaswami.
\newblock CN: Verifying systems C code with separation-logic refinement types.
\newblock In \emph{Proc. POPL}, 2023.

\bibitem{cogent} S.~Amani et al.
\newblock Cogent: Verifying high-assurance file system implementations.
\newblock In \emph{Proc. ASPLOS}, 2016.

\bibitem{seldma} BlueRock Security.
\newblock Verified ZynqMP DMA driver in concurrent separation logic.
\newblock seL4 Summit (presentation), 2025.

\bibitem{embassy} The \texttt{embassy} project and \texttt{embedded-dma}.
\newblock \url{https://embassy.dev}.

\bibitem{ravenscar} A.~Burns and A.~Wellings.
\newblock \emph{Concurrent and Real-Time Programming in Ada} (the Ravenscar
profile).
\newblock Cambridge University Press, 2007.

\bibitem{an4839} STMicroelectronics.
\newblock AN4839: Level 1 cache on STM32F7 and STM32H7 series.
\newblock Application note.

\bibitem{easydma} Nordic Semiconductor.
\newblock nRF52832 Product Specification --- EasyDMA (``EasyDMA is not able to
access the Flash'').

\bibitem{zephyrdma} Zephyr Project.
\newblock DMA and data cache coherency on Arm M7 devices.
\newblock GitHub issue zephyrproject-rtos/zephyr \#36471.

\bibitem{stm32eth} stm32-rs/stm32-eth.
\newblock Document possible problems when placing buffers on a cacheable region.
\newblock GitHub issue stm32-rs/stm32-eth \#16.

\bibitem{nrfhal} nrf-rs/nrf-hal.
\newblock EasyDMA doesn't read from Flash.
\newblock GitHub issue nrf-rs/nrf-hal \#37.

\end{thebibliography}

\end{document}
